\thispagestyle{myheadings}
\addcontentsline{toc}{chapter}{\protect\numberline {} 参考文献}

\begin{thebibliography}{参考文献}

  % \bibitem{Colby}
  % Latha, S. Colby and Dirk VanGucht,
  % ``A Grammar Model for Database'',
  % {\it TECHNICAL REPORT,} NO.282,
  % June 1989.

  % \bibitem{Gonnet}
  % Gaston, H. Gonnet and Frank Wm. Tompa,
  % ``Mind Your Grammar: a New Approach to Modelling Text'',
  % {\it Proceedings of the 13th VLDB Conference,}Brighton,
  % pp. 339--346, 1987.

  % \bibitem{Dzenan}
  % Dzenan RIDJANOVIC and Micheal L. BRODIE,
  % ``DEFINING DATABASE DYNAMICS WITH ATTRIBUTE GRAMMARS'',
  % {\it INFORMATION PROCESSING LETTERS,}vol. 14, No. 3,
  % pp. 132--138, May 1982.

  \bibitem{inf}
  Robert E. Kraut, Robert S. Fish, Robert W. Root, and Barbara L. Chalfonte,
  ``Informal Communication in Organizations: Form, Function, and Technology'',
  In {\sl Human Reactions to Technology: Claremont Symposium on Applied Social Psychology,}
  S. Oskamp and S. Spacapan (eds.),
  Sage Publications,
  1990,
  pp.145--199.

  \bibitem{tel}
  国土交通省,
  ``令和5年度 テレワーク人口実態調査 -調査結果(概要)-'',
  国土交通省,
  2024,
  \url{https://www.mlit.go.jp/report/press/content/001733057.pdf},
  (参照 2026-1-14).

  \bibitem{fle}
  厚生労働省,
  ``労働時間制度等に関するアンケート調査結果について（速報値）'',
  厚生労働省　労働基準局,
  2024,
  \url{https://www.mhlw.go.jp/content/11201250/001194510.pdf},
  (参照 2026-1-14).


  \bibitem{BTIW2023}
  情報処理学会，
  ``Behavior Transformation by IoT International Workshop (BTIW2023)'',
  情報処理学会IoT行動変容学研究グループ,
  2023,
  \url{https://www.sig-bti.jp/event/btiw.html},
  (参照 2026-1-14).

  \bibitem{nadge}
  池本忠弘.
  ``ナッジと行動変容―科学に基づく新しい政策アプローチ―''
  {\sl 保健医療科学,}
  Vol.73,
  No.4,
  2024,
  pp.256--264.

  \bibitem{nadge:body}
  竹林正樹，甲斐裕子，江口泰正，西村司，山口大輔，福田洋.
  ``わかっていてもなかなか実践しない相手をどう動かす？—身体活動・運動促進へのナッジ—''
  {\sl 第29 回日本健康教育学会学術大会,}
  Vol.30,
  No.1,
  2022,
  pp.73--78.

  \bibitem{feedback:sit}
  道城裕貴，松見淳子.
  ``通常学級において「めあて\&フィードバックカード」による目標設定とフィードバックが着席行動に及ぼす効果''
  {\sl 行動分析学研究,}
  Vol.20,
  No.2,
  2007,
  pp.118--128.

  \bibitem{game:health}
  Hiroaki Uechi, Nobusuke Tan, Yuichiro Honda,
  ``Effects of gamification-based intervention for promoting health behaviors'',
  {\sl The Journal of Physical Fitness and Sports Medicine,}
  Vol.7,
  No.3,
  2018,
  pp.185--192.

  \bibitem{CGM}
  Maxine E. Whelan, Francesca Denton, Claire L. A. Bourne, Andrew P. Kingsnorth, Lauren B. Sherar, Mark W. Orme, and Dale W. Esliger,
  ``A digital lifestyle behaviour change intervention for the prevention of type 2 diabetes: a qualitative study exploring intuitive engagement with real-time glucose and physical activity feedback'',
  In {\sl BMC Public Health,}
  Vol.21,
  No.1,
  2021,
  pp.130.

  \bibitem{cafe}
  中野利彦，亀和田慧太，杉戸準，永岡良章，小倉加奈代，西本一志.
  ``Traveling Cafe：分散型オフィス環境におけるコミュニケーション促進支援システム''
  {\sl インタラクション2006論文集,}
  Vol.4,
  2006,
  pp.227--228.

  \bibitem{irori}
  松原孝志，臼杵正郎，杉山 公造，西本 一志.
  ``言い訳オブジェクトとサイバー囲炉裏 : 共有インフォーマル空間におけるコミュニケーションを触発するメディアの提案''
  {\sl 情報処理学会論文誌,}
  Vol.44,
  2003,
  pp.3174--3187.

  \bibitem{awareness}
  Paul Dourish, and Victoria Bellotti,
  ``Awareness and Coordination in Shared Workspaces'',
  In {\sl Proceedings of the ACM Conference on Computer-Supported Cooperative Work (CSCW),}
  ACM,
  1992,
  pp.107--114.

  \bibitem{gotogether}
  Minghao Cai, and Jiro Tanaka,
  ``Go together: providing nonverbal awareness cues to enhance co-located sensation in remote communication'',
  In {\sl Proceedings of the ACM Conference on Human Factors in Computing Systems (CHI),}
  ACM,
  2019,
  pp.1--12.

  \bibitem{yohou}
  田中優斗, 福島拓, 吉野孝.
  ``Docoitter：未来の在室情報を予報する在室管理システム''
  {\sl 情報処理学会論文誌,}
  Vol.54,
  No.9,
  2013,
  pp.2265--2275.

  \bibitem{hayashi}
  林航平，梶克彦.
  ``未来の在室状況予測と共通の話題に基づく来訪促進システムの提案''
  {\sl マルチメディア，分散協調とモバイル シンポジウム 2025 論文集,}
  Vol.2025,% 書き方怪しいので一応先生に確認
  2025,
  pp.1552--1561.

  \bibitem{ikusei}
  濱田もえ，仲西渉，安尾萌，松下光範.
  ``育成ゲーム型アプリケーションを用いた研究室滞在の習慣化に向けた一検討''
  {\sl 情報処理学会第79回全国大会講演論文集,}
  No.4,
  2017,
  pp.829--830.

  \bibitem{huneas}
  松田完，西本 一志.
  ``HuNeAS : 大規模組織内での偶発的な出会いを利用した情報共有の促進とヒューマンネットワーク活性化支援の試み''
  {\sl 情報処理学会論文誌,}
  Vol.43,
  2002,
  pp.3571--3581.

  \bibitem{concierge}
  森田篤史，山下邦弘，國藤進.
  ``インタレスト コンシェルジェ"待ち状況"における共通興味を案内情報提供サービスシステム''
  {\sl インタラクション2003論文集,}
  2003,
  pp.189--190.

  \bibitem{kenkou}
  田木真和，玉木悠，森川富昭，松久宗英，森口博基.
  ``NFC通信歩数計を活用した健康データの可視化による生活習慣の行動変容''
  {\sl 医療情報学,}
  Vol.34,
  2014,
  pp.281--291.

  \bibitem{kion}
  深谷恭平，村上暁信.
  ``可視化された温熱環境情報が屋外空間の利用に与える影響に関する研究''
  {\sl 都市計画論文集,}
  Vol.58,
  2023,
  pp.1439--1445.

  \bibitem{zinro}
  粟生将之，梶克彦.
  ``生体情報共有に基づく人狼ゲームの拡張に関する研究''
  In {\sl Proceedings of the Workshop on Informatics 2025 (WiNF 2025),}
  2025,
  pp.59--62.

  \bibitem{kancolle}
  中島瑞希，梶克彦.
  ``環境コレクション''
  In {\sl Proceedings of the Workshop on Informatics 2025 (WiNF 2025),}
  2025,
  pp.17--20.


  \bibitem{fingerprint:WiFi}
  Moustafa Youssef, and Ashok Agrawala,
  ``The Horus WLAN Location Determination System''
  In {\sl Proceedings of the 3rd International Conference on Mobile Systems, Applications, and Services (MobiSys '05),}
  ACM Press,
  2005,
  pp.205--218.

  \bibitem{fingerprint:BLE}
  Ramsey Faragher, and Robert Harle,
  ``Location Fingerprinting With Bluetooth Low Energy Beacons''
  {\sl IEEE Journal on Selected Areas in Communications,}
  Vol.33,
  No.11,
  2015,
  pp.2418--2428.

  \bibitem{threeposition:WiFi}
  Santiago Mazuelas, Alfonso Bahillo, Ruben M. Lorenzo, Patricia Fernandez, Francisco A. Lago, Eduardo Garcia, Juan Blas, and Evaristo J. Abril,
  ``Robust Indoor Positioning Provided by Real-Time RSSI Values in Unmodified WLAN Networks,''
  {\sl IEEE Journal of Selected Topics in Signal Processing,}
  Vol.3,
  No.5,
  2009,
  pp.821--831.

  \bibitem{threeposition:BLE}
  Aitor De Blas, and Diego López-de-Ipiña,
  ``Improving trilateration for indoors localization using BLE beacons,''
  In {\sl Proceedings of the 2017 2nd International Multidisciplinary Conference on Computer and Energy Science (SpliTech 2017),}
  2017,
  pp.1--6.


  \bibitem{PDR:position}
  上坂大輔，村松茂樹，岩本健嗣，横山浩之.
  ``手に保持されたセンサを用いた歩行者向けデッドレコニング手法の提案''
  {\sl 情報処理学会論文誌,}
  Vol.52,
  No.2,
  2011,
  pp.558--570.

  \bibitem{PDR:direction}
  星尚志，藤井雅弘，羽多野裕之，伊藤篤，渡辺裕.
  ``スマートフォンを用いた歩行者デッドレコニングのための進行方向推定に関する研究''
  {\sl 情報処理学会論文誌,}
  Vol.57,
  No.1,
  2016,
  pp.25--33.

  \bibitem{BLE:Proximity}
  Sylvia T. Kouyoumdjieva, and Gunnar Karlsson,
  ``Experimental Evaluation of Precision of a Proximity-based Indoor Positioning System,''
  in {\sl Proc.\ of IEEE International Conference on Indoor Positioning and Indoor Navigation (IPIN),}
  2019.


  \bibitem{mapmatching}
  吉見駿，村尾和哉，望月祐洋，西尾信彦.
  ``マップマッチングを用いたPDR軌跡補正''
  {\sl 情報処理学会研究報告書,}
  No.20,
  2014,
  pp.1--8.

  \bibitem{landmark}
  Samuel House, Sean Connell, Ian Milligan, Daniel Austin, Tamara L. Hayes, and Patrick Chiang,
  ``Indoor localization using pedestrian dead reckoning updated with RFID-based fiducials,''
  In {\sl Proceedings of the 33rd Annual International Conference of the IEEE Engineering in Medicine and Biology Society (EMBS 2011),}
  2011,
  pp. 7598--7601.


  \bibitem{GoogleMap}
  Google，
  ``Google Maps'',
  \url{https://www.google.com/maps}，
  (参照 2026-01-14).

  \bibitem{PokemonGo}
  Scopely，
  ``Pokémon GO'',
  \url{https://pokemongo.com/}，
  (参照 2026-01-14).


  \bibitem{layout}
  Daisuke Kawamura, Tsuyoshi Setoguchi.
  ``Walking-trajectory-based survey analysis of pedestrians’ actual usages of open spaces in a large-scale commercial facility''
  {\sl Journal of Asian Architecture and Building Engineering,}
  Vol.12,
  Article 33,
  2025.

  \bibitem{pattern}
  鈴木直彦，平澤宏祐，田中健一，小林貴訓，佐藤洋一，藤野陽三.
  ``人物動線データ群における逸脱行動人物検出及び行動パターン分類''
  {\sl 電子情報通信学会論文誌. D, 情報・システム,}
  Vol.91,
  No.6,
  2008,
  pp.1550--1560.


  \bibitem{simulation}
  沖拓弥，小川芳樹.
  ``大地震時を想定した広域避難シミュレーション軌跡のクラスタリングと避難行動予測への応用''
  {\sl 人工知能学会全国大会論文集,}
  第34回全国大会,
  2020.

  \bibitem{light}
  山下俊樹，三木光範，中原蒼太，間博人，高谷友貴.
  ``執務者が BLE ビーコンを携帯するビーコン携帯型知的照明システム''
  {\sl 情報科学技術フォーラム(FIT),}
  Vol.16,
  No.4,
  2017,
  pp.345--346.

  \bibitem{airControl}
  進藤宏行，鶴見隆太，鈴木義康，早川慶朗，長廣剛，磯崎日出雄，竹林英樹.
  ``屋外への開放部を持つ空間における人流・気流センサを用いた空調制御手法の開発・実証 （第１報）開発・実証システムの概要と実証1年目の結果''
  {\sl 空気調和・衛生工学会大会 学術講演論文集,}
  Vol.2019.9,
  2019,
  pp.173--176.


  \bibitem{staywatch}
  Fumitaka Naruse, Katsuhiko Kaji,
  ``Estimation of Person Existence in Room Using BLE Beacon and Its Platform''
  {\sl In Proceedings of the 2018 International Conference on Mobile Computing, Applications, and Services (MobiCASE),}
  Vol.240,
  2018,
  pp.251--257.


  \bibitem{togawa}
  戸川浩汰，梶克彦.
  ``滞在推定システムにおけるプライバシを考慮した携帯型BLEビーコンの設計と実装''
  {\sl 研究報告コンシューマ・デバイス\&システム(CDS),}
  Vol.2025-CDS-43,
  No.6,
  2025,
  pp.1--8.


  \bibitem{ICcard}
  須田憲人，鯉沼辰弥，鈴木崇，熊澤弘之.
  ``非接触 IC カードを用いた出席登録システム''
  {\sl 電子情報通信学会技術研究報告,}
  Vol.111,
  No.141,
  2011,
  pp.65--70.


  \bibitem{camera}
  株式会社ビットキー，
  ``workhub 顔認証ソリューション'',
  \url{https://www.bitlock.workhub.site/product/face-recognition}，
  (参照 2026-01-14).


  \bibitem{voice}
  石澤鴻弥，岩井将行.
  ``発話とその認識エリアに基づく在室管理システムの提案''
  {\sl 研究報告ユビキタス コンピューティングシステム(UBI),}
  Vol.57,
  No.19,
  2018,
  pp.1--6.


  \bibitem{WiFi}
  東昭太朗，曽超.
  ``Web 在室管理システムにおける無線 LAN を用いた入退室処理の自動化への試み''
  {\sl 第 81 回全国大会講演論文集,}
  Vol.1,
  2019,
  pp.115--116.

  \bibitem{BLE:EnvironmentBeacon}
  山下陸，西山勇毅，小松寛弥，川原圭博.
  ``BLE ビーコンを用いた屋内位置推定システムの設計と実装''
  {\sl 	研究報告ユビキタスコンピューティングシステム(UBI),}
  Vol.2020-UBI-68,
  No.8,
  2020,
  pp.1--7.

  \bibitem{BLE:UserBeacon}
  Yashar Kiarashi, Soheil Saghafi, Barun Das, Chaitra Hegde, Venkata Siva Krishna Madala, ArjunSinh Nakum, Ratan Singh, Robert Tweedy, Matthew Doiron, Amy D. Rodriguez, Allan I. Levey, Gari D. Clifford, and Hyeokhyen Kwon,
  ``Graph Trilateration for Indoor Localization in Sparsely Distributed Edge Computing Devices in Complex Environments Using Bluetooth Technology,''
  {\sl Sensors,}
  Vol. 23,
  No. 23,
  2023.


  \bibitem{RaspberryPi}
  Raspberry Pi Foundation，
  ``Raspberry Pi'',
  \url{https://www.raspberrypi.com/products/}，
  (参照 2026-01-14).


  \bibitem{opendata}
  国土交通省，
  ``公共交通オープンデータセンター'',
  国土交通省，
  \url{https://www.odpt.org/}，
  (参照 2026-1-14).

  \bibitem{navi}
  株式会社ナビタイムジャパン，
  ``NAVITIME API'',
  株式会社ナビタイムジャパン，
  \url{https://api-sdk.navitime.co.jp/}，
  (参照 2026-1-14).

  \bibitem{eki}
  株式会社ヴァル研究所，
  ``駅すぱあとWebサービス'',
  株式会社ヴァル研究所，
  \url{https://api.ekispert.jp/}，
  (参照 2026-1-14).


  \bibitem{ondoku}
  音読さん，
  ``音読さん（オンラインテキスト読み上げサービス）'',
  音読さん，
  \url{https://ondoku3.com/ja/}，
  (参照 2026-1-14).


  \bibitem{slack}
  Slack Technologies，
  ``Slack'',
  \url{https://slack.com/}，
  (参照 2026-01-14).



  % \bibitem{ifo}
  % Abiteboul, S. and Hull, R.
  % ``IFO: A formal semantic database model''
  % {\sl ACM Transactions on Database Systems  12,} 4,
  % December 1987,
  % pp.525-565.
  %
  %\bibitem{er}
  %Chen, P. P. 
  %``The entitiy-relationship model-toward a unified view of data''
  %{\sl ACM Transactions on Database Systems 1,} 1,
  %March 1976,
  %pp.9-36.
  %
  % \bibitem{sdm}
  % Hammer, M. and McLeod, D.
  % ``Database discription with SDM: A semantic database model''
  % {\sl ACM Transactions on Database Systems 6,} 3,
  % September 1981,
  % pp.351-386.
  %
  %\bibitem{fdm}
  %Shipman, D. W.
  %``The functioal data model and the language DAPLEX''
  %{\sl ACM Transactions on Database Systems 6,} 1,
  %March 1981,
  %pp.140-173.
  %
  %\bibitem{o2}
  %Bancilhon, F., et al.
  %``The design and implementation of O$_2$, an object-oriented database system''
  %{\sl Processings of 2nd International Workshop on Object-Oriented Database Systems}
  %1988,
  %pp.1-22.
  %
  %\bibitem{gem}
  %Copeland, G., and Maier, D.
  %``Making Smalltalk a database system''
  %{\sl Processings of the Annual SIGMOD Conference,}
  %June 1984.
  %
  %\bibitem{iris}
  %Fishman, D. H., et al. 
  %``Iris: An object oriented database management system''
  %{\sl ACM Transactions on Database Systems 5,} 1,
  %January 1987,
  %pp.48-69.
  %
  %\bibitem{orion}
  %Kim, W.,Manerjee, J., Chou, H. T., Garza, J. F., and Woelk, D.
  %``Composite object support in an object-oriented database system''
  %{\sl Proceedings of OOPSLA,}
  %1987,
  %pp.118-125.

  % \bibitem{Hull}
  % Hull, R. and Yap, C. K.
  % ``The Format model: A theory of database organization'',
  % {\it JACM 31,}3,
  % pp. 518--537, 1984.

  % \bibitem{jap1}
  % 増永良文,
  % ``次世代データベースシステムとしてのオブジェクト指向データベースシステム'',
  % 情報処理,Vol. 32, No. 5,
  % pp. 490--499, May 1991.

  % \bibitem{jap2}
  % 田中克己,
  % ``オブジェクト指向データベースの基礎概念''
  % 情報処理,Vol. 32, No. 5,
  % May 1991,
  % pp. 500-513.

  % \bibitem{omt}
  % J.ランボー M.プラハ W.プレメラニ F.エディ W.ローレンセン,
  % ``OBJECT-ORIENTED MODELING AND DESIGN'',
  % トッパン, 1992.

  % \bibitem{model}
  % 米澤明憲,
  % ``モデル化と表現''
  % 岩波書店(1992).

  % \bibitem{prolog}
  % 中島秀之,
  % ``Prolog''
  % 産業図書(1983).

\end{thebibliography}
% Local Variables: 
% mode: japanese-LaTeX
% TeX-master: "root"
% End: 
